\honest{Original Seeing}{Eliezer Yudkowsky}{2007}

Robert Pirsig put this well, so I copy what he said. ``I don't know if this story is based on reality or not, but either way, it's true.'' \nb{Pirsig's Author's Note says the book is ``based on actual occurrences'' and ``must be regarded in its essence as fact,'' though much was changed. If the story were invented, it could not also be the evidence for its lesson, and the post gives no other. Two days later, in ``The Logical Fallacy of Generalization from Fictional Evidence,'' Yudkowsky asked whether an issue does not ``have to be argued in its own right.''}

The rest of the post is Pirsig's. A writing teacher in Bozeman has students who cannot think of anything to say. One wants to write five hundred words about the United States. He tells her to narrow it down to Bozeman, then to its main street, and she still has nothing, and comes back in tears. Furious, he says, ``You're not looking!'' For every fact there is an infinity of hypotheses, and the more you look the more you see. He tells her: ``Narrow it down to the front of one building on the main street of Bozeman. The Opera House. Start with the upper left-hand brick.''

She comes back with five thousand words. ``I sat in the hamburger stand across the street,'' she says, ``and started writing about the first brick, and the second brick, and then by the third brick it all started to come and I couldn't stop.''

The teacher concludes that she had been trying to repeat what she had already heard, and that narrowing the subject to one brick forced her to see for herself. \nb{That is close to the idea of ``Cached Thoughts,'' posted three days earlier; the post leaves the connection unstated. In the book, the passage goes on: the teacher had whole classes write for an hour about the back of his thumb, and then a coin, with no complaints about having nothing to say. The post stops before this.}
